% © 2026 Ing. David Jaroš — CC BY-NC-ND 4.0
%
% This work is licensed under a Creative Commons Attribution-NonCommercial-NoDerivatives
% 4.0 International License (CC BY-NC-ND 4.0).
%
% License History: Earlier drafts (up to v0.3) were released under CC BY 4.0.
% From v0.4 onward, all material is released under CC BY-NC-ND 4.0 to protect
% the integrity of the theoretical work during ongoing academic development.

\documentclass[12pt]{article}
\usepackage{amsmath,amssymb,amsthm}
\usepackage{geometry}
\usepackage{hyperref}
\usepackage{tcolorbox}
\usepackage{xcolor}
\geometry{margin=1in}

% Theorem-like environments
\newtheorem{definition}{Definition}
\newtheorem{proposition}{Proposition}
\newtheorem{theorem}{Theorem}
\newtheorem{remark}{Remark}
\newtheorem{conjecture}{Conjecture}

\title{ST-3: Stability of Imaginary Sector Configurations in UBT\\[0.5em]
\large\textsc{Speculative Extension — Invisibility Track}}
\author{David Jaroš}
\date{2026-03-03}

\begin{document}
\maketitle

\begin{tcolorbox}[colback=yellow!10, title=Status]
\textbf{SPECULATIVE EXTENSION} — This content is not part of
the canonical UBT framework.  It is an exploratory investigation
into the stability of imaginary-sector field configurations and has not
been validated against the canonical field equations.
\end{tcolorbox}

\vspace{1em}
\noindent\textit{Output of Sub-task ST-3 from the Invisibility / Imaginary Sector
investigation track.  See also:
\texttt{st1\_imaginary\_object\_definition.tex},
\texttt{st2\_sector\_interaction.tex},
\texttt{st4\_ctc\_invisibility\_connection.tex},
\texttt{st5\_falsifiability.md}.}

%----------------------------------------------------------------------
\section*{Abstract}
%----------------------------------------------------------------------
We analyse whether an imaginary-sector field configuration
$\Theta_I = i\,\widetilde\Theta_I$ is stable within UBT.
Three questions are addressed:
(i) does the UBT Hamiltonian $H[\Theta_I]$ have a lower bound?
(ii) can $\Theta_I$ decay into $\Theta_R$ via some process?
(iii) is there a conserved charge that forbids sector mixing?
We find that: (i) the Hamiltonian density of $\Theta_I$ equals
that of $\widetilde\Theta_I$ (same sign, same lower bound);
(ii) sector-mixing transitions are forbidden at leading order
by the discrete symmetry $\Theta \to i\Theta$; and
(iii) a candidate conserved charge (the imaginary-sector
particle number) is identified.  All conclusions are \textbf{speculative}.

%----------------------------------------------------------------------
\section{Hamiltonian of the Imaginary Sector}
%----------------------------------------------------------------------

From ST-2, the energy-momentum tensor of $\Theta_I = i\,\widetilde\Theta_I$
(with $\widetilde\Theta_I$ real-valued) equals that of $\widetilde\Theta_I$:
\begin{equation}
  T_{\mu\nu}[\Theta_I]
  \;=\; T_{\mu\nu}[\widetilde\Theta_I].
  \label{eq:tmunu_id}
\end{equation}
The energy density (Hamiltonian density) in flat space is
\begin{equation}
  \mathcal{H}[\Theta_I]
  \;=\; T_{00}[\Theta_I]
  \;=\; |\dot{\widetilde\Theta}_I|^2
       + |\nabla\widetilde\Theta_I|^2
       + V(\widetilde\Theta_I^\dagger\widetilde\Theta_I).
  \label{eq:hamiltonian_density}
\end{equation}
The first two terms are manifestly non-negative.  The third term
$V(\widetilde\Theta_I^\dagger\widetilde\Theta_I)$ is bounded below
if and only if the UBT potential $V$ is bounded below — which is a
requirement of the canonical UBT framework (stable vacuum).

\begin{proposition}[Lower bound on $H[\Theta_I]$]
\label{prop:lower_bound}
Under the assumption that $V \geq V_\mathrm{min} > -\infty$
(canonical UBT stability condition), the Hamiltonian of the imaginary
sector is bounded below:
\begin{equation}
  H[\Theta_I]
  \;=\; \int d^3x\,\mathcal{H}[\Theta_I]
  \;\geq\; V_\mathrm{min} \cdot \mathrm{Vol},
\end{equation}
where $\mathrm{Vol}$ is the spatial volume.  The imaginary sector
does not introduce any new instability: its energy functional has
the same lower bound as the real sector.
\end{proposition}

%----------------------------------------------------------------------
\section{Sector-Mixing Transitions}
%----------------------------------------------------------------------

A sector-mixing transition would convert an imaginary-sector quantum
into a real-sector quantum:
\begin{equation}
  \Theta_I\text{-quanta} \;\longrightarrow\; \Theta_R\text{-quanta} + \text{(radiation)}.
\end{equation}

\subsection{Leading-Order Amplitude}

The interaction vertex that couples $\Theta_R$ and $\Theta_I$ comes
from the cross-term in the action (from ST-2, Proposition~2):
\begin{equation}
  S_\mathrm{cross}
  \;\sim\; \int d^4x\,\sqrt{-g}\,
  \left[\operatorname{Im}(\Theta_R^\dagger\widetilde\Theta_I)\right]^2
  V''(\Theta_R^\dagger\Theta_R).
\end{equation}
This is a quartic interaction of the form
$(\Theta_R \cdot \Theta_I)^2$, coupling four field quanta.
For a process $\Theta_I \to \Theta_R + X$ the amplitude is at
least second order in the coupling, and the amplitude vanishes unless
there are at least two imaginary-sector quanta in the initial state.

\subsection{Symmetry Argument}

Consider the global phase transformation
\begin{equation}
  \Theta \;\mapsto\; e^{i\theta}\,\Theta, \qquad \theta \in \mathbb{R}.
  \label{eq:u1_sym}
\end{equation}
If the potential $V$ depends only on $|\Theta|^2 = \Theta^\dagger\Theta$,
this is a $U(1)$ symmetry of the action.  Under this symmetry, the
imaginary sector and real sector mix under $\theta = \pi/2$:
\begin{equation}
  e^{i\pi/2}\,\Theta_R \;=\; i\,\Theta_R \;=\; \Theta_I.
\end{equation}
This means $\Theta_R$ and $\Theta_I = i\,\Theta_R$ lie in the
\emph{same} $U(1)$ orbit.  In a theory with exact $U(1)$ symmetry,
the distinction between $\Theta_R$ and $\Theta_I$ is a matter of
phase convention and cannot correspond to distinct sectors with
different observable properties.

\paragraph{Breaking of $U(1)$ by electromagnetic coupling.}
Standard model matter couples to $\Theta_R$ (the real sector) via
gauge fields.  This coupling breaks the $U(1)$ symmetry
$\Theta \to e^{i\theta}\Theta$ explicitly for $\theta \neq 0, \pi$,
creating a genuine distinction between $\Theta_R$ and $\Theta_I$.
After this breaking, the $U(1)$ is reduced to a discrete $\mathbb{Z}_2$:
\begin{equation}
  \mathbb{Z}_2:\;\Theta \;\mapsto\; -\Theta.
  \label{eq:z2}
\end{equation}
The imaginary sector $\Theta_I \to i\Theta_R$ is related to the
real sector by a $\mathbb{Z}_4$ rotation
$(\Theta \to i\Theta)$, which is broken by the electromagnetic
coupling.

\begin{conjecture}[Sector mixing rate]
\label{conj:mixing_rate}
In a UBT theory where the electromagnetic coupling explicitly breaks
$U(1) \to \mathbb{Z}_4$ (or lower), sector-mixing transitions
$\Theta_I \to \Theta_R$ are \emph{allowed but suppressed} by the
mixing-breaking scale.  The decay rate is estimated as
\begin{equation}
  \Gamma(\Theta_I \to \Theta_R)
  \;\sim\; \frac{g_\mathrm{mix}^2\,m_I^5}{192\pi^3\,m_\mathrm{Pl}^4},
\end{equation}
where $g_\mathrm{mix}$ is the sector-mixing coupling constant,
$m_I$ is the imaginary-sector particle mass, and $m_\mathrm{Pl}$
is the Planck mass.  For $g_\mathrm{mix} \ll 1$ and $m_I \ll m_\mathrm{Pl}$
this rate is extremely small and the imaginary sector is
effectively stable on cosmological timescales.
\end{conjecture}

%----------------------------------------------------------------------
\section{Conserved Charge and Stability}
%----------------------------------------------------------------------

\subsection{Imaginary-Sector Particle Number}

Under the residual $\mathbb{Z}_4$ symmetry $\Theta \to i\Theta$, a
Noether-like charge can be defined by treating the complex phase of
$\Theta$ as a $U(1)$ degree of freedom before gauge fixing.  The
associated Noether current is:
\begin{equation}
  J^\mu \;=\; i(\Theta^\dagger\partial^\mu\Theta - \partial^\mu\Theta^\dagger\,\Theta),
  \label{eq:noether_current}
\end{equation}
with conservation $\partial_\mu J^\mu = 0$ following from the $U(1)$
equations of motion.

\paragraph{Decomposition into real and imaginary sectors.}
Inserting $\Theta = \Theta_R + \Theta_I$:
\begin{equation}
  J^\mu
  \;=\; J^\mu[\Theta_R] + J^\mu[\Theta_I]
       + i\!\left(
           \Theta_R^\dagger\partial^\mu\Theta_I
         - \partial^\mu\Theta_R^\dagger\,\Theta_I
         + \mathrm{h.c.}
         \right).
\end{equation}
The cross-term in $J^\mu$ is purely imaginary (by the same argument
as in ST-2), so the real conserved charge decomposes as
\begin{equation}
  Q \;=\; \int d^3x\,J^0
  \;=\; Q[\Theta_R] + Q[\Theta_I].
\end{equation}
The imaginary-sector particle number
$Q_I = Q[\Theta_I]$ is separately conserved at leading order.

\begin{theorem}[Conservation of imaginary-sector charge]
\label{thm:charge_conservation}
If the UBT action respects the global $U(1)$ symmetry \eqref{eq:u1_sym}
before electromagnetic coupling, then the imaginary-sector particle
number $Q_I$ is conserved at leading order in the coupling constant.
Sector-mixing processes that change $Q_I$ are suppressed by powers
of the electromagnetic coupling.  The imaginary sector is therefore
\emph{perturbatively stable}.
\end{theorem}

\subsection{Stability Classification}

Two regimes arise depending on the electromagnetic coupling strength:

\begin{enumerate}
  \item \textbf{Exact $U(1)$ symmetry} (no electromagnetic coupling,
    or coupling that respects the full $U(1)$): $Q_I$ is exactly
    conserved, imaginary-sector objects are absolutely stable.
    Invisibility is permanent.

  \item \textbf{Broken $U(1)$} (standard model electromagnetic
    coupling): $Q_I$ is only approximately conserved.
    Imaginary-sector objects can decay, but the lifetime is
    $\tau \sim m_\mathrm{Pl}^4 / (g_\mathrm{mix}^2 m_I^5)$, which
    can be much longer than the age of the universe for
    $m_I \ll m_\mathrm{Pl}$.  Invisibility is long-lived but transient.
\end{enumerate}

%----------------------------------------------------------------------
\section{Vacuum Stability and Sector Separation}
%----------------------------------------------------------------------

A further stability question is whether the imaginary-sector vacuum
$\langle\Theta_I\rangle \neq 0$ is energetically preferred over the
real-sector vacuum $\langle\Theta_R\rangle \neq 0$.

In the canonical UBT framework the vacuum is determined by the
minimum of $V(\Theta^\dagger\Theta)$.  Since $\Theta^\dagger\Theta$
is the same for $\Theta = \Theta_R$ and $\Theta = i\,\Theta_R =
\Theta_I$ (because $|i\Theta_R|^2 = |\Theta_R|^2$), the vacuum
energy is degenerate between the real and imaginary sectors.

\begin{remark}[Degenerate vacuum]
The UBT potential $V(\Theta^\dagger\Theta)$ has a continuous set
of degenerate minima related by $\Theta \to e^{i\theta}\Theta$.
The distinction between the real-sector vacuum and the imaginary-sector
vacuum is a choice of $U(1)$ phase — a gauge convention, not a
physical distinction, unless the $U(1)$ is broken by interactions.
This supports the picture in which invisible (imaginary-sector)
matter is a valid, degenerate vacuum sector of UBT.
\end{remark}

%----------------------------------------------------------------------
\section{Summary}
%----------------------------------------------------------------------
\begin{enumerate}
  \item The Hamiltonian of the imaginary sector is bounded below
    (Proposition~\ref{prop:lower_bound}), so imaginary-sector
    configurations do not introduce new instabilities.
  \item Sector-mixing transitions are suppressed by the explicit
    breaking of the $U(1)$ symmetry by electromagnetic coupling.
    The imaginary sector is perturbatively stable
    (Theorem~\ref{thm:charge_conservation}).
  \item A conserved charge $Q_I$ (imaginary-sector particle number)
    is identified; it forbids sector mixing at leading order.
  \item The vacuum degeneracy between real and imaginary sectors
    supports the existence of imaginary-sector matter as a valid
    physical state of UBT.
  \item Invisibility is \emph{stable} in the absence of
    electromagnetic coupling, and \emph{long-lived} (potentially
    cosmological) when electromagnetic coupling is present.
\end{enumerate}

\medskip
\noindent\textbf{All statements in this document are speculative.
No claim constitutes an established result of the canonical UBT framework.}

%----------------------------------------------------------------------
\bibliographystyle{plain}
\begin{thebibliography}{9}
\bibitem{UBT_core}
  D.~Jaroš,
  \emph{Unified Biquaternion Theory: core field equations},
  Repository document \texttt{consolidation\_project/ubt\_core\_main.tex} (2025).
\bibitem{ST1}
  D.~Jaroš,
  \emph{ST-1: Imaginary Sector Objects in UBT},
  Repository document
  \texttt{speculative\_extensions/invisibility/st1\_imaginary\_object\_definition.tex}
  (2026).
\bibitem{ST2}
  D.~Jaroš,
  \emph{ST-2: Interaction Between Real and Imaginary Sectors in UBT},
  Repository document
  \texttt{speculative\_extensions/invisibility/st2\_sector\_interaction.tex}
  (2026).
\bibitem{Kolb1990}
  E.~W.~Kolb and M.~S.~Turner,
  \emph{The Early Universe},
  Addison-Wesley (1990).
\end{thebibliography}

\end{document}
